% =============================================================================
% Lista de abreviaturas (páginas preliminares). La lista de símbolos matemáticos está en el
% anexo de notación (\listadesimbolos, definida en notacion.tex).
% =============================================================================
\chapter*{Lista de abreviaturas}
\addcontentsline{toc}{chapter}{Lista de abreviaturas}
\markboth{Lista de abreviaturas}{Lista de abreviaturas}
{\small
\begin{longtable}{@{}p{0.2\textwidth}p{0.74\textwidth}@{}}
\toprule
Abreviatura & Significado \\
\midrule
\endhead
\bottomrule
\endfoot
ADF & Prueba de Dickey-Fuller aumentada (raíz unitaria). \\
AE & Autoencoder. \\
\texttt{aede} & Función de EPCA con datación del paquete \texttt{AedeClahfac} del Dr.~Camiz. \\
ARI & Índice de Rand ajustado. \\
BCRP & Banco Central de Reserva del Perú. \\
BVL & Bolsa de Valores de Lima. \\
CH & Índice de Caliński y Harabasz. \\
\texttt{clahfac} & Función de HFC del paquete \texttt{AedeClahfac}. \\
D1--D19 & Decisiones metodológicas registradas en \texttt{decisiones.qmd} (\cref{tab:decisiones}). \\
DS & Serie integrada (\emph{difference stationary}), según la regla D5. \\
EHFC & HFC evolutiva (objetivo condicional OC1). \\
EMBIG & Diferencial soberano del Perú en el índice de bonos de mercados emergentes. \\
EPCA & Análisis evolutivo de componentes principales. \\
FMI & Fondo Monetario Internacional. \\
HAC & Estimador de varianza consistente ante heterocedasticidad y autocorrelación. \\
HFC & Clasificación Factorial Jerárquica (\emph{Hierarchical Factor Classification}). \\
HT1--HT4 & Hipótesis de trabajo (\cref{sec:hipotesis}). \\
IPC & Índice de precios al consumidor. \\
KMO & Medida de adecuación muestral de Kaiser-Meyer-Olkin. \\
KPSS & Prueba de estacionariedad alrededor de una tendencia (\cref{sec:met-diagnostico}). \\
LSTM & Red recurrente de memoria de largo y corto plazo (\emph{long short-term memory}). \\
M1, M3 & Agregados monetarios: dinero y liquidez total. \\
NBER & National Bureau of Economic Research (Estados Unidos). \\
OC1--OC4 & Objetivos exploratorios condicionales (\cref{sec:oe}). \\
OE1--OE12 & Objetivos específicos (\cref{sec:oe}). \\
OMS & Organización Mundial de la Salud. \\
PBI & Producto bruto interno. \\
PCA & Análisis de componentes principales. \\
PCA-MM\ventanaMediaMovil{} & Datación con los \emph{scores} del PCA suavizados con media móvil causal (\cref{def:datacion}). \\
RV & Coeficiente RV de Robert y Escoufier. \\
SBS & Superintendencia de Banca, Seguros y AFP. \\
STATIS & Método de comparación de tablas múltiples (\emph{Structuration des Tableaux À Trois Indices de la Statistique}). \\
TAMN & Tasa activa promedio en moneda nacional. \\
TIPMN & Tasa pasiva promedio en moneda nacional. \\
TS & Serie estacionaria en tendencia (\emph{trend stationary}), según la regla D5. \\
USTR & Oficina del Representante Comercial de los Estados Unidos. \\
\end{longtable}
}
